\documentclass[12pt,a4paper]{article}
\usepackage[utf8]{vietnam}
\usepackage{amsmath}
\usepackage{amssymb}
\usepackage{listings}
\usepackage{xcolor}
\usepackage{geometry}

\geometry{
    left=2.5cm,
    right=2.5cm,
    top=2.5cm,
    bottom=2.5cm
}

\lstset{
    language=C,
    basicstyle=\ttfamily\small,
    keywordstyle=\color{blue}\bfseries,
    commentstyle=\color{green!60!black},
    stringstyle=\color{red},
    numbers=left,
    numberstyle=\tiny\color{gray},
    frame=single,
    breaklines=true,
    showstringspaces=false,
    tabsize=4
}

\title{\textbf{Bài tập ôn tập lập trình C cơ bản - HW1}}
\author{}
\date{}

\begin{document}

\maketitle

\section{Bài 1: Xâu ký tự và cấp phát bộ nhớ động}

\subsection{Đề bài}
Xây dựng và chạy minh họa hàm:
\begin{verbatim}
char* subStr(char* s1, int offset, int number);
\end{verbatim}

Hàm tách xâu con từ xâu s1 bắt đầu từ ký tự tại chỉ số offset (tính từ 0) và có độ dài number ký tự.

\subsection{Yêu cầu}
\begin{itemize}
    \item Kiểm tra tính hợp lệ của các đối số
    \item Nếu number lớn hơn độ dài phần còn lại của xâu s1 từ vị trí offset, trả về phần còn lại của s1 từ vị trí offset
    \item Không dùng mảng ký tự tĩnh
    \item Không sử dụng các hàm xử lý xâu có sẵn trong thư viện string.h
\end{itemize}

\subsection{Input/Output}
\textbf{Input:}
\begin{itemize}
    \item Dòng thứ nhất: Hai số nguyên offset và number
    \item Dòng thứ hai: Xâu str (tối đa 80 ký tự)
\end{itemize}

\textbf{Ví dụ:}
\begin{verbatim}
6 3
hust.edu.vn
\end{verbatim}

\textbf{Output:}
\begin{verbatim}
-edu-
\end{verbatim}

\subsection{Thuật toán}
\begin{enumerate}
    \item Kiểm tra tính hợp lệ của offset và number:
    \begin{itemize}
        \item Nếu offset < 0 hoặc offset >= strlen(s1): trả về NULL
        \item Nếu number <= 0: trả về NULL
    \end{itemize}
    \item Tính độ dài thực tế của xâu con:
    \begin{itemize}
        \item Nếu offset + number > strlen(s1): length = strlen(s1) - offset
        \item Ngược lại: length = number
    \end{itemize}
    \item Cấp phát bộ nhớ động cho xâu con (length + 1 ký tự cho '\textbackslash0')
    \item Copy từng ký tự từ s1[offset] đến xâu con
    \item Thêm ký tự '\textbackslash0' vào cuối xâu con
    \item Trả về xâu con
\end{enumerate}

\subsection{Code C}
\begin{lstlisting}
#include <stdio.h>
#include <stdlib.h>

int my_strlen(char* s) {
    int len = 0;
    while (s[len] != '\0') {
        len++;
    }
    return len;
}

char* subStr(char* s1, int offset, int number) {
    // Kiểm tra tính hợp lệ
    if (s1 == NULL || offset < 0 || number <= 0) {
        return NULL;
    }
    
    int len = my_strlen(s1);
    if (offset >= len) {
        return NULL;
    }
    
    // Tính độ dài thực tế
    int actual_length;
    if (offset + number > len) {
        actual_length = len - offset;
    } else {
        actual_length = number;
    }
    
    // Cấp phát bộ nhớ động
    char* result = (char*)malloc((actual_length + 1) * sizeof(char));
    if (result == NULL) {
        return NULL;
    }
    
    // Copy ký tự
    for (int i = 0; i < actual_length; i++) {
        result[i] = s1[offset + i];
    }
    result[actual_length] = '\0';
    
    return result;
}

int main() {
    int offset, number;
    char str[81];
    
    printf("Nhap offset va number: ");
    scanf("%d %d", &offset, &number);
    getchar(); // Xóa ký tự newline
    
    printf("Nhap xau str: ");
    fgets(str, 81, stdin);
    
    // Xóa ký tự newline nếu có
    int len = my_strlen(str);
    if (str[len - 1] == '\n') {
        str[len - 1] = '\0';
    }
    
    char* sub = subStr(str, offset, number);
    if (sub != NULL) {
        printf("-%s-\n", sub);
        free(sub);
    } else {
        printf("Invalid input\n");
    }
    
    return 0;
}
\end{lstlisting}

\newpage

\section{Bài 2: Cấu trúc}

\subsection{Đề bài}
Cho hai cấu trúc:
\begin{verbatim}
typedef struct point {
    double x;
    double y;
} point_t;

typedef struct circle {
    point_t center;
    double radius;
} circle_t;
\end{verbatim}

Viết hàm \texttt{is\_in\_circle} trả về 1 nếu điểm p nằm trong đường tròn c.

\subsection{Hàm cần viết}
\begin{verbatim}
int is_in_circle(point_t *p, circle_t *c);
\end{verbatim}

\subsection{Input/Output}
\textbf{Input:}
\begin{itemize}
    \item Dòng thứ nhất: 3 số thực biểu diễn tọa độ x, y và bán kính của tâm đường tròn
    \item Dòng thứ hai: 2 số thực biểu diễn tọa độ x, y của điểm p
\end{itemize}

\textbf{Ví dụ:}
\begin{verbatim}
0 0 5.5
2.3 -3.1
\end{verbatim}

\textbf{Output:}
\begin{itemize}
    \item YES nếu p nằm trong đường tròn
    \item NO nếu p không nằm trong đường tròn (p nằm trên hoặc nằm ngoài đường tròn)
\end{itemize}

\subsection{Thuật toán}
Điểm p nằm trong đường tròn c khi khoảng cách từ p đến tâm đường tròn nhỏ hơn bán kính:
\begin{equation}
\sqrt{(p.x - c.center.x)^2 + (p.y - c.center.y)^2} < c.radius
\end{equation}

Để tránh tính căn bậc hai, ta so sánh bình phương:
\begin{equation}
(p.x - c.center.x)^2 + (p.y - c.center.y)^2 < c.radius^2
\end{equation}

\subsection{Code C}
\begin{lstlisting}
#include <stdio.h>
#include <math.h>

typedef struct point {
    double x;
    double y;
} point_t;

typedef struct circle {
    point_t center;
    double radius;
} circle_t;

int is_in_circle(point_t *p, circle_t *c) {
    double dx = p->x - c->center.x;
    double dy = p->y - c->center.y;
    double distance_squared = dx * dx + dy * dy;
    double radius_squared = c->radius * c->radius;
    
    if (distance_squared < radius_squared) {
        return 1; // Nằm trong đường tròn
    } else {
        return 0; // Nằm trên hoặc ngoài đường tròn
    }
}

int main() {
    circle_t c;
    point_t p;
    
    printf("Nhap toa do tam (x y) va ban kinh: ");
    scanf("%lf %lf %lf", &c.center.x, &c.center.y, &c.radius);
    
    printf("Nhap toa do diem p (x y): ");
    scanf("%lf %lf", &p.x, &p.y);
    
    if (is_in_circle(&p, &c)) {
        printf("YES\n");
    } else {
        printf("NO\n");
    }
    
    return 0;
}
\end{lstlisting}

\end{document}
